\subsection{Adjacent-page supplement: Kontsevich and Morozov}
\label{supp:1985-adjacent}
These two pages are preserved as supplementary evidence only. Their records
667--668 lie outside the requested core records 641--666; they are excluded
from the year's core talk and occurrence counts and receive ADJ-prefixed
identifiers. The manuscript heading on P001258 reads \(24/\mathrm{II}/85\),
whereas catalogue record 667 gives 24 February 1987. P001259 reads
24 February 1986, whereas record 668 also gives 24 February 1987.
Neither conflict is silently repaired. Kontsevich's matching paper was
submitted 3 December 1986 and published in 1987; Morozov explicitly reports
Belavin--Knizhnik and Beilinson rather than claiming new results. The technical
assessment below is independent of resolving these dates.

\subsubsection{ADJ-667.1: A curve with a point and a formal parameter is not unmarked Teichmueller space}
\noindent\textbf{Source.} P001258; M. L. Kontsevich.\par
\noindent\textbf{Notebook.} Kontsevich's page calls \(M_g\) Teichm\"uller space and writes \(\widetilde M_g=(C,p,x)\), where \(x\) is a coordinate at \(p\).\par

\noindent\textbf{Precise statement.} For \(g\ge2\), let \(\widehat{\mathcal M}_{g,1}\) parameterize smooth connected projective complex curves \(C\) of genus \(g\), a point \(p\), and a formal parameter \(x\in\widehat{\mathfrak m}_p\) with nonzero first derivative. It is the inverse limit of the finite-jet moduli spaces. Forgetting \(x\) gives pointed moduli \(\mathcal M_{g,1}\); forgetting both \(p,x\) gives \(\mathcal M_g\). A topological marking would instead place these constructions over Teichm\"uller space.\par

\noindent\textbf{Reconstruction and limits.} The triple records no topological marking. We use a moduli stack or a local smooth cover rather than silently identify it with Teichm\"uller space. A formal parameter need not be a convergent analytic coordinate on a fixed disk.\par

\noindent\textbf{Literature fate.} \citet[\S1]{kontsevich1987virasoro} explicitly constructs these finite-jet moduli and their inverse limit. The paper was submitted 3 December 1986 and published in issue 2 of 1987.\par

\noindent\textbf{Present status.} As of November 2026: The objects are identifiable standard deformation-theoretic data (E). The notebook's two uses of the word Teichm\"uller require this distinction.\par

\noindent\textbf{Assessment.} Class E; literal correctness: correct with missing hypotheses; publication match: exact; verification: read the result; development: none; confidence: high.\par

\subsubsection{ADJ-667.2: Formal local vector fields and the curve-dependent global subalgebra}
\noindent\textbf{Source.} P001258; M. L. Kontsevich.\par
\noindent\textbf{Notebook.} The page defines \(\mathfrak g=\mathbb C((x))\partial_x\), \(\mathfrak k=\mathbb C[[x]]\partial_x\), and \(\gamma\) as Laurent expansions of vector fields on \(C\) with a pole only at \(p\).\par

\noindent\textbf{Precise statement.} Set \(\gamma_{C,p,x}=H^0(C,T_C(*p))\subset\mathbb C((x))\partial_x\) by Laurent expansion. Its elements are algebraic vector fields on \(C\setminus\{p\}\), meromorphic on \(C\) with finite pole order at \(p\). The local subalgebra \(\mathfrak k\) consists of fields regular on the formal disk, while \(\mathfrak k_0=x\mathbb C[[x]]\partial_x\) consists of fields fixing its marked origin.\par

\noindent\textbf{Reconstruction and limits.} The space \(\gamma\) depends on the curve and parameter; \(\mathfrak k\) does not. The marginal phrase \ru{всю C!!} cannot mean holomorphic extension through \(p\): for \(g\ge2\), \(H^0(C,T_C)=0\), whereas meromorphic fields with allowed pole order are essential here. Editorial check of the marginal whole-curve interpretation \ours.\par

\noindent\textbf{Literature fate.} These subalgebras and the global punctured-curve condition are defined in \citet[\S1]{kontsevich1987virasoro}.\par

\noindent\textbf{Present status.} As of November 2026: The finite-pole global subalgebra is exactly identified (E). The ambiguous marginal instruction is retained as provenance rather than imposed as a false extra condition.\par

\noindent\textbf{Assessment.} Class E; literal correctness: correct with missing hypotheses; publication match: exact; verification: read the result; development: annotated; confidence: high.\par

\subsubsection{ADJ-667.3: Virasoro uniformization: the infinitesimal action and stabilizer}
\noindent\textbf{Source.} P001258; M. L. Kontsevich.\par
\noindent\textbf{Notebook.} Theorem1 says \(\mathfrak g\) acts on \(\widetilde M_g\), with stabilizer \(\gamma\), so it is locally the homogeneous space \(\mathfrak g/\gamma\).\par

\noindent\textbf{Precise statement.} For the moduli of smooth genus-\(g\ge2\) pointed curves with a formal parameter, the Lie algebra \(\operatorname{Der}\mathbb C((x))\) acts infinitesimally. At \((C,p,x)\), evaluation of that action is surjective and its kernel is \(H^0(C,T_C(*p))\). Thus \(T_{(C,p,x)}\widehat{\mathcal M}_{g,1}\cong\mathbb C((x))\partial_x/\gamma_{C,p,x}\), in the formal inverse-limit sense.\par

\noindent\textbf{Reconstruction and limits.} This is an infinitesimal uniformization theorem. It does not assert the existence of an ordinary global Lie group, a single fixed subgroup for all curves, or a global homogeneous-space presentation. The completed formal setting and genus hypothesis are supplied.\par

\noindent\textbf{Literature fate.} Exact match to \citet[\S1, proposition and preceding tangent formula]{kontsevich1987virasoro}. The paper credits independent results of Beilinson and Schechtman. Publication in 1987 follows either the manuscript's 1985 heading or the catalogue's February 1987 date; the latter conflict is not resolved by the mathematics.\par

\noindent\textbf{Present status.} As of November 2026: Published soon after by the speaker (A), with the precise infinitesimal meaning verified in the primary paper.\par

\noindent\textbf{Assessment.} Class A; literal correctness: correct with missing hypotheses; publication match: exact; verification: read the result; development: none; confidence: high.\par

\subsubsection{ADJ-667.4: A meromorphic lift of a one-parameter deformation}
\noindent\textbf{Source.} P001258; M. L. Kontsevich.\par
\noindent\textbf{Notebook.} For a one-parameter family the notebook chooses a connection with poles only at the section \(p\), locally \(\nabla_t=\partial_t+F(x,t)\partial_x\).\par

\noindent\textbf{Precise statement.} Locally on a smooth one-dimensional base, a family of smooth projective curves with section admits a lift of the base tangent vector to a relative meromorphic vector field after allowing a sufficiently high finite pole order along that section. In a formal trivialization it has the written form. Two such lifts differ by a relative vertical meromorphic field; at a fixed fiber their difference belongs to \(\gamma_{C,p,x}\).\par

\noindent\textbf{Reconstruction and limits.} Allowing a sufficiently large pole divisor kills the relevant first cohomology by Serre vanishing, giving local existence. This is a noncanonical local lift, not a globally chosen holomorphic connection. Its ambiguity explains the quotient by \(\gamma\). The notebook neither fixes a pole order nor proves uniqueness. Editorial local-lift argument from vanishing of the first cohomology after a sufficiently large pole divisor \ours.\par

\noindent\textbf{Literature fate.} This is the deformation-theoretic interpretation of \citet[\S1]{kontsevich1987virasoro}; the article states the equivalent Kodaira--Spencer and Laurent-field description rather than this exact local connection formula.\par

\noindent\textbf{Present status.} As of November 2026: The construction is explained within the established infinitesimal uniformization. It is E as deformation-theoretic machinery, with a partial publication match rather than a claim of a separately stated theorem.\par

\noindent\textbf{Assessment.} Class E; literal correctness: correct with missing hypotheses; publication match: partial; verification: read the result; development: annotated; confidence: high.\par

\subsubsection{ADJ-667.5: The erroneous g/k tangent quotient and the two forgetting maps}
\noindent\textbf{Source.} P001258; M. L. Kontsevich.\par
\noindent\textbf{Notebook.} One line says \(T\widetilde M_g=\mathfrak g/\mathfrak k\), but the following display divides instead by fields extending to \(C\setminus\{p\}\).\par

\noindent\textbf{Precise statement.} The tangent space with full formal parameter is \(\mathfrak g/\gamma\). The tangent space of unpointed moduli is \(\mathfrak g/(\gamma+\mathfrak k)\cong H^1(C,T_C)\), of complex dimension \(3g-3\). The tangent space of pointed moduli is \(\mathfrak g/(\gamma+\mathfrak k_0)\cong H^1(C,T_C(-p))\), of dimension \(3g-2\). The intermediate quotient \(\mathfrak g/\mathfrak k\) is merely the space of formal principal parts; it is not any of these complete tangent spaces.\par

\noindent\textbf{Reconstruction and limits.} These formulas follow from the two-open-set formal gluing calculation: changes of the trivialization on the punctured curve contribute \(\gamma\), while those on the disk contribute \(\mathfrak k\) or \(\mathfrak k_0\) depending on whether its center is remembered. They also distinguish \(\operatorname{Der}\mathbb C[[x]]\) from the Lie algebra \(x\mathbb C[[x]]\partial_x\) of formal coordinate changes fixing the point. Editorial formal-gluing derivation of the two forgetting-map tangent quotients \ours; the full-parameter uniformization theorem is due to Kontsevich.\par

\noindent\textbf{Literature fate.} \citet[\S1]{kontsevich1987virasoro} gives the full-parameter tangent formula and notes that regular formal fields act transitively along the fibers of forgetting the point and parameter. The two finite-dimensional cohomology quotients are the explicit editorial reconstruction.\par

\noindent\textbf{Present status.} As of November 2026: The first written quotient is false (E correction); the next display contains the intended curve-dependent stabilizer. No open deformation problem remains.\par

\noindent\textbf{Assessment.} Class E; literal correctness: false as written; publication match: exact; verification: read the result; development: annotated; confidence: high.\par

\subsubsection{ADJ-667.6: The central extension lifts to determinant-of-cohomology lines}
\noindent\textbf{Source.} P001258; M. L. Kontsevich.\par
\noindent\textbf{Notebook.} The final lines say that the Virasoro central extension acts in the determinant bundle \(d_k\) pulled back from \(M_g\) to \(\widetilde M_g\).\par

\noindent\textbf{Precise statement.} For an integer \(k\), let \(\lambda_k(C)=\det H^0(C,K_C^k)\otimes\det H^1(C,K_C^k)^{-1}\). The infinitesimal Laurent-vector-field action lifts to a Virasoro central extension acting by first-order differential operators on the pullback of \(\lambda_k\) to curves with a formal parameter. Relative to the determinant extension for \(k=0\) (equivalently \(k=1\)), its cocycle class is multiplied by \(6k^2-6k+1\). The numerical central charge changes with central-generator and dual-line conventions.\par

\noindent\textbf{Reconstruction and limits.} The notebook does not define \(d_k\); identifying it with \(\lambda_k\) is the exact source-based reconstruction. A Lie algebra action on sections with nontrivial central action is not an ordinary action of the center-free algebra on the base. No particular scalar central charge is assigned without declaring the cocycle normalization.\par

\noindent\textbf{Literature fate.} \citet[\S\S2--3]{kontsevich1987virasoro} constructs the determinant extension on the Laurent-space Grassmannian and pulls it back to the curve moduli. Section3 supplies the coefficient and the corresponding cancellation corollary.\par

\noindent\textbf{Present status.} As of November 2026: Published by the speaker in 1987 (A). The determinant-line action is verified directly; no unrelated later Kontsevich-model claim is attached to the page.\par

\noindent\textbf{Assessment.} Class A; literal correctness: correct with missing hypotheses; publication match: exact; verification: read the result; development: none; confidence: high.\par

\subsubsection{ADJ-668.1: The critical bosonic measure and the exponent13}
\noindent\textbf{Source.} P001259; Morozov.\par
\noindent\textbf{Notebook.} Morozov, under ``Consequences of the Belavin--Knizhnik theorem'', summarizes Beilinson's talk and writes a sum over genera with measure \(|d\mu_p|^2/(\det\operatorname{Im}T)^{13}\), multiplied by a vertex factor \(K\{V\}\).\par

\noindent\textbf{Precise statement.} For a smooth compact Riemann surface of fixed genus \(p\ge2\), the gauge-fixed critical bosonic vacuum measure can be expressed locally as \(|\mu_p|^2/(\det\operatorname{Im}\tau)^{13}\), where \(\tau\) is the period matrix in a symplectic marking and \(\mu_p\) is the corresponding holomorphic Mumford form, a section of the determinant-line combination \(\lambda_2\otimes\lambda_1^{-13}\). Insertions require the associated normalized matter/ghost correlation factor and their integration data.\par

\noindent\textbf{Reconstruction and limits.} This is a fixed-genus determinant identity with specified zero modes. The infinite genus sum and integration near degenerating curves are not proved convergent. Genus 0 and 1 have automorphisms and additional special zero-mode conventions, so they are not covered by silently repeating the same formula. The scan confirms the exponent13.\par

\noindent\textbf{Literature fate.} \citet[\S\S2--5, especially (2.19), (3.4), and (5.1)]{belavinknizhnik1986complex} proves the determinant-line/holomorphy statement and discusses degeneration poles. The paper cites its own Landau Institute 1985 preprint; the full article was submitted 18 February 1986 and published August 1986.\par

\noindent\textbf{Present status.} As of November 2026: Exposition of already circulating work (E), explicitly credited in the notebook. The recorded theorem does not justify a finite total string partition function.\par

\noindent\textbf{Assessment.} Class E; literal correctness: correct with missing hypotheses; publication match: exact; verification: read the result; development: none; confidence: high.\par

\subsubsection{ADJ-668.2: World sheets, topology, and replacing external circles by punctures}
\noindent\textbf{Source.} P001259; Morozov.\par
\noindent\textbf{Notebook.} Two sketches show a propagating loop and a surface capped at two points. The text says that interactions are accounted for by different topologies, and that it is technically convenient to begin and end at points.\par

\noindent\textbf{Precise statement.} In the perturbative world-sheet description of closed strings, the genus labels the loop order, while punctures and their assigned states describe external insertions. Turning an external cylindrical end into a puncture requires a local coordinate and a state/vertex prescription. This is a formal sewing description of amplitudes, not a topological theorem that arbitrary prescribed boundary-loop data equal a point delta function.\par

\noindent\textbf{Reconstruction and limits.} The caption abbreviation is retained without a forced expansion. The diagrams specify the physical construction but not a measure, sewing theorem, or convergence statement. The point insertion model below is one chosen Fourier representation.\par

\noindent\textbf{Literature fate.} The earlier surface-sum and insertion formulas are in \citet[eqs.~(1)--(3),(26)--(31)]{polyakov1981quantum}; the fixed-genus consequences appear in \citet[\S6]{belavinknizhnik1986complex}.\par

\noindent\textbf{Present status.} As of November 2026: Known formal perturbative interpretation (E); no new precise topological or analytic question is posed.\par

\noindent\textbf{Assessment.} Class E; literal correctness: not determinable; publication match: thematic only; verification: read the result; development: none; confidence: high.\par

\subsubsection{ADJ-668.3: Fourier-transforming target point constraints into exponential vertices}
\noindent\textbf{Source.} P001259; Morozov.\par
\noindent\textbf{Notebook.} The page first calls \(V_\alpha\) delta functions, then writes \(V_\alpha(p)=\int e^{ip_\alpha\cdot X(\xi)}d^2\xi\).\par

\noindent\textbf{Precise statement.} For a fixed smooth map, Fourier transformation of \(\delta^{(d)}(X(\xi)-y)\) in the target point \(y\), with convention \(\int e^{ip\cdot y}\delta(X-y)dy\), gives \(e^{ip\cdot X(\xi)}\). In the quantum free-boson theory this exponential needs normal ordering. An integrated closed-string vertex must also be a coordinate-invariant world-sheet density; in conformal coordinates the usual on-shell scalar vertex has conformal weights \((1,1)\) before integration against \(d^2z\).\par

\noindent\textbf{Reconstruction and limits.} The elementary Fourier identity is exact for a smooth map. Its promotion to an integrated physical vertex requires a normalization of the action, a momentum/mass-shell condition, regularization, and the ghost prescription. None is specified by the bare coordinate integral. Editorial Fourier identity in the stated sign convention \oursComputed; the quantum vertex construction is prior literature.\par

\noindent\textbf{Literature fate.} \citet[eqs.~(26)--(31)]{polyakov1981quantum} gives the vertex construction; \citet[\S6, (6.21)--(6.26)]{belavinknizhnik1986complex} computes the resulting tachyon factors with momentum conservation and mass-shell conditions.\par

\noindent\textbf{Present status.} As of November 2026: Established construction (E), with the distinction between a classical Fourier identity and a renormalized quantum vertex retained.\par

\noindent\textbf{Assessment.} Class E; literal correctness: correct with missing hypotheses; publication match: exact; verification: read the result; development: annotated; confidence: high.\par

\subsubsection{ADJ-668.4: A target delta derivative is not a world-sheet derivative vertex}
\noindent\textbf{Source.} P001259; Morozov.\par
\noindent\textbf{Notebook.} The sentence about derivatives of delta functions is followed by \(\int\partial X^\mu e^{ip\cdot X}\,d^2\xi\).\par

\noindent\textbf{Precise statement.} Fourier transformation of a target-coordinate derivative \(\partial_{y^\mu}\delta(X-y)\) gives a factor \(-ip_\mu e^{ip\cdot X}\) under the preceding convention. It does not give \(\partial X^\mu\). World-sheet derivatives of the free field produce oscillator or descendant insertions. For a closed-string tensor vertex a typical full density includes both \(\partial X^\mu\) and \(\bar\partial X^\nu\), appropriate polarization and on-shell conditions; the isolated chiral factor has an uncontracted world-sheet index.\par

\noindent\textbf{Reconstruction and limits.} The notebook does not say which delta variable is differentiated. If it means the same target point whose Fourier transform was just taken, the indicated implication is false. If it means a different state construction, that is missing data, not a justification for silently adding the antiholomorphic factor. Editorial target-delta derivative calculation \oursComputed.\par

\noindent\textbf{Literature fate.} The distinction is an elementary Fourier calculation, checked against the exponential-insertion setup in \citet[eqs.~(26)--(31)]{polyakov1981quantum}.\par

\noindent\textbf{Present status.} As of November 2026: E with a faithful mathematical gap; no general new vertex-operator theorem is inferred from the incomplete line.\par

\noindent\textbf{Assessment.} Class E; literal correctness: false as written; publication match: partial; verification: read the result; development: annotated; confidence: high.\par

\subsubsection{ADJ-668.5: Gaussian integration with sources, determinant powers, and the zero mode}
\noindent\textbf{Source.} P001259; Morozov.\par
\noindent\textbf{Notebook.} Integrating over \(X\) is said to give \((\det\Delta_0)^{-13}\) times products of exponentials \(p_\alpha p_\beta G_{\mathrm{reg}}(\xi_\alpha,\xi_\beta)\). The preceding action has a positive exponent.\par

\noindent\textbf{Precise statement.} Fix a closed connected surface and a nonnegative scalar Laplacian \(\Delta\). For 26 real mean-zero fields with action \(S=\frac12\langle X,\Delta X\rangle\), a finite-mode Gaussian integral with sources \(\sum_a p_a\delta_{z_a}\) equals a normalization times \((\det{}'\Delta)^{-13}\exp[-\frac12\sum_{a,b}p_a\cdot p_b\,G(z_a,z_b)]\), before diagonal regularization, where \(G=\Delta^{-1}\) on mean-zero functions. Integrating the constant target mode separately gives \((2\pi)^{26}\delta^{(26)}(\sum_a p_a)\). The infinite-dimensional determinant requires zeta regularization.\par

\noindent\textbf{Reconstruction and limits.} The sign and factor 1/2 are stated for this explicit action and covariance convention; another convention for \(G\) changes their appearance. A positive Euclidean exponential of positive energy is divergent. The prime, momentum-conservation delta, and contractions of target indices are missing from the notebook's shorthand. Derivative vertices require additional differentiated Green-function contractions. Editorial finite-mode Gaussian completion of the square and zero-mode check \oursComputed; the string-theory application is the cited published calculation.\par

\noindent\textbf{Literature fate.} \citet[\S6, (6.21)--(6.24)]{belavinknizhnik1986complex} performs this Gaussian source calculation and explicitly uses momentum conservation. The finite-mode completion-of-the-square identity provides the normalization check.\par

\noindent\textbf{Present status.} As of November 2026: Established Gaussian calculation (E). The result read is a regularized fixed-surface building block, not a rigorous value for the total genus sum.\par

\noindent\textbf{Assessment.} Class E; literal correctness: false as written; publication match: exact; verification: read the result; development: annotated; confidence: high.\par

\subsubsection{ADJ-668.6: The Green-function diagonal and Wick ordering}
\noindent\textbf{Source.} P001259; Morozov.\par
\noindent\textbf{Notebook.} The final line defines \(G_{\mathrm{reg}}\) as a Green function smoothed when the two points coincide.\par

\noindent\textbf{Precise statement.} The two-dimensional scalar covariance has a logarithmic diagonal singularity. For a mollified mean-zero Gaussian field \(X_\varepsilon\), define a point exponential at finite cutoff by \(:e^{ip\cdot X_\varepsilon(z)}:=\exp(ip\cdot X_\varepsilon(z)+\frac12|p|^2\mathbb E[X_\varepsilon^1(z)^2])\). At pairwise distinct points, its normalized correlation has limit \(\exp[-\sum_{a<b}p_a\cdot p_bG(z_a,z_b)]\). Integrated collisions, local metric factors, and alternative diagonal finite parts require additional prescriptions.\par

\noindent\textbf{Reconstruction and limits.} Here the components have identical scalar covariance and are independent. This formula identifies one precise diagonal subtraction. Merely replacing the diagonal by an arbitrary smooth value does not select the physical normalization or prove integrability when insertion points collide. Editorial Wick-subtraction and off-diagonal contraction calculation \oursComputed; normal ordering itself is established machinery.\par

\noindent\textbf{Literature fate.} \citet[\S6, between (6.24) and (6.26)]{belavinknizhnik1986complex} explicitly regularizes the diagonal and then imposes the mass-shell condition; \citet[\S3]{dotsenkofateev1984conformal} supplies the corresponding normal-ordered free-boson contraction framework.\par

\noindent\textbf{Present status.} As of November 2026: Known renormalized Gaussian building block (E). No unqualified convergence theorem is extracted from the notebook's word ``smoothed''.\par

\noindent\textbf{Assessment.} Class E; literal correctness: correct with missing hypotheses; publication match: partial; verification: read the result; development: annotated; confidence: high.\par

